% TLDR: Distinguish semantic validity controls, reportable outcomes, and private data.
\section{Semantic Controls and Reporting}
\label{sec:controls}
Agreement with experts using a faulty reference can reproduce the same error rather than measure coverage correctly.
We therefore accompany the random audit with controls whose semantic relationship is checked independently wherever possible.
These controls do not receive the random audit's sampling interpretation, and their counts and selection process are reported separately.

Removing a necessary obligation should reduce coverage, provided no remaining statement already implies it.
Narrowing a precondition should reduce coverage when it excludes inputs required by the task.
Equivalent formulations should receive similar ratings, while adding irrelevant or vacuous statements should not increase meaningful coverage.
Changing proof bodies without changing statements should preserve a statement-coverage rating.
An incomplete or false reference should trigger discrepancy reporting where the evidence supports it.
The exact tolerances and control counts are pilot decisions, not assumed measured properties of the judge.

For equivalence controls, compilation of two programs is insufficient.
The relevant implications must preserve domains, assumptions, and obligation weights.
For negative controls, a counterexample should establish that the removed obligation is not already entailed.
AI can assist in constructing these certificates, but checker acceptance and the assumptions of the certificate must be recorded.
Unresolved cases remain unresolved rather than being converted into negative examples by a timeout.

The reference-anchoring study is a separate controlled comparison within this research program.
It varies reference representation and proof visibility while holding the intended obligations and candidate artifact fixed.
Its labels, estimand, and pilot counts must not be substituted for those of the main conditional-rating study.
Likewise, the human-rating study does not change an existing benchmark's prover-based scoring implementation.

\subsection{What the final report must contain}
The primary table will report RPS, paired baseline differences, and confidence intervals for all 896 audit tasks, the 128 seen tasks with fresh outputs, the 768 held-out tasks, and each solver configuration.
The four excluded-solver fits have their own transfer table and denominators.
Distribution-error diagnostics on the dense subset, means, signed bias, absolute error, Pearson correlation, Spearman correlation, and Kendall correlation supplement these results.
Human--human agreement is computed on matched artifacts, and the prediction mean is not confused with the full probability vector.

A separate repeatability panel reports variation across repeated probability vectors and expected scores, the number of full judge executions, and the aggregation policy.
A stable but wrong judge can pass a repeatability check, while a correct prediction of human disagreement can intentionally be broad.
The paper must keep these statements separate.
Missing and failed calls appear with their original denominator; because RPS is bounded, worst- and best-case completion bounds can show whether missingness could reverse a conclusion.

Public evidence consists of the protocol, configurations, input manifests where distributable, analysis code, and authorized aggregate results.
Participant identities, contact information, consent, individual rating sheets, and sealed test ratings remain in restricted storage.
Opaque analysis identifiers preserve the assignment structure without exposing identities.
Any eventual release of individual ratings requires the corresponding consent and release decision.
The historical import preserves immutable aggregate receipts and hashes, with post hoc calculations labeled separately.
